\section{Fluid benchmarks and instance optimality}
\label{sec:common-fluid-envelopes}

The randomized instance-optimality results all have the same leading-order
explanation.  Fix an input vector and privately shuffle its values among the
public job labels.  Uniformly over predictable prefixes, the cumulative
number of encountered jobs of each type differs by only $o(n)$ from the
prefix length times its full-input proportion.  Thus, if a fraction $t$ of
the jobs has been tested, the revealed mass of type $p_i$ is approximately
$tD[p](\{p_i\})$.  As long as a nonnegligible fraction remains, the normalized
composition of the untouched jobs is consequently still close to $D[p]$.

The fluid model makes this approximation exact.  It replaces $n$ jobs by one
unit of divisible job mass and divides machine work by $n$.  Testing mass
$0.01$, for example, reveals exactly a $0.01$ copy of the input distribution.
This removes the random fluctuations and the rounding of job counts, leaving
a deterministic scheduling problem for the leading $n^2$ term of total
completion time.

Despite the name, a fluid schedule will not be a policy that makes a new
decision at every real time.  The distribution below has finitely many job
types, and a schedule is a finite sequence of divisible blocks.  Within one
block, work and completions accrue proportionally.  This linear traversal is
the limit of dividing the block into many finite batches; the later
subsections justify the corresponding approximation for actual jobs.  The
fluid optimization developed first contains no sampling or limiting
argument: it solves the finite-dimensional problem exactly when the
distribution is known.  Only afterward do we transfer that solution to
shuffled finite inputs.

The same problem covers three models.  We prescribe a total tested fraction
$q$ and let $v$ be the work per job in a completion mode that does not reveal
the processing time in advance.  Obligatory testing has $q=1$, revealing
optimization has $v=u$, and blind execution has
$v=\int p\,dD(p)$.  Blind optimization is different because its preliminary
operation does not reveal the processing time, so it is not represented by
the testing blocks below.

\subsection{Fluid blocks and schedules}
\label{sec:fluid-schedules}

We first define the fluid problem operationally.  This makes clear both what
an arbitrary fluid schedule may do and why its state remains finite.  Let
\begin{equation}
 D=\sum_{i=0}^m d_i\delta_{p_i},
 \qquad 0=p_0<p_1<\cdots<p_m,
 \qquad \sum_{i=0}^m d_i=1,
 \label{eq:fluid-finite-distribution}
\end{equation}
where $d_0$ may be zero.  Fix $q\in[0,1]$ and $v\ge0$.

A fluid schedule is a finite sequence of the following operations.

\begin{description}[leftmargin=*,style=nextline]
 \item[Testing block of mass $z$ with selector $g$.]
 The vector $g=(g_1,\ldots,g_m)\in[0,1]^m$ specifies which fraction of each
 positive type is processed directly after being revealed.  The block tests
 mass $z$, completes all revealed zero jobs and the selected positive jobs,
 and places the remaining revealed positive jobs in deferred queues.  It
 uses work
 \[
   z\left(1+\sum_{i=1}^m p_i d_i g_i\right)
 \]
 and completes mass
 \[
   z\left(d_0+\sum_{i=1}^m d_i g_i\right).
 \]
 It adds $zd_i(1-g_i)$ to the deferred queue of type $i$.  The total mass of
 all testing blocks may not exceed $q$.

 \item[Deferred-processing block of type $i$.]
 Processing mass $z$ from the currently available deferred queue of type
 $i$ uses work $p_i z$ and completes mass $z$.

 \item[Alternative-completion block.]
 Completing mass $z$ in the other mode uses work $vz$ and completes mass
 $z$.  The total mass of these blocks may not exceed $1-q$.
\end{description}

The schedule is complete if it tests total mass $q$, uses the other mode on
mass $1-q$, and eventually empties every deferred queue.  A positive-work
block is traversed linearly; a zero-work block completes its mass at once.
Consequently every finite fluid schedule has a piecewise-linear completion
curve $C(x)$, the mass completed after $x$ units of work.

The area under the unfinished mass is its total completion cost:
\begin{equation}
 \Area(C)\defeq\int_0^\infty(1-C(x))\,dx.
 \label{eq:completion-curve-area}
\end{equation}
Indeed, if a mass element $z$ completes after work $T(z)$, the elementary
layer-cake identity gives
\[
 \int T(z)\,dz
 =\int_0^\infty \operatorname{mass}\{z:T(z)>x\}\,dx
 =\Area(C).
\]
Thus continuous work is only a convenient coordinate for drawing the
completion curve; the schedule itself makes finitely many block decisions.

\subsection{The best testing block}
\label{sec:fluid-testing-block}

We next determine which revealed positive types should be completed inside a
testing block.  For the finite fluid model, all integrals below are finite
sums under \cref{eq:fluid-finite-distribution}.  The density lemma is no more
difficult for an arbitrary finite-mean distribution, so we state it in that
form for later use.

For a selector $g$, let
\begin{equation}
 a_D(g)\defeq D(\{0\})+\int_{(0,\infty)}g(p)\,dD(p),
 \qquad
 w_D(g)\defeq 1+\int_{(0,\infty)}p g(p)\,dD(p).
 \label{eq:common-module-aw}
\end{equation}
These are the completed mass and work when one unit of mass is tested.
Their ratio $a_D(g)/w_D(g)$ is the block's \emph{completion density}.

For example, suppose half the mass has length zero and half has length one.
Tests alone complete mass $1/2$ using work $1$.  Processing the length-one
outcomes as well completes the full mass using work $3/2$, improving the
completion density from $1/2$ to $2/3$.  For longer positive outcomes the
same decision can reduce the density.  At the optimum, a positive type is
included precisely when its length is below the work per completion of the
resulting block.

\begin{lemma}[Maximum-density testing block]
\label{lem:maximum-density-module}
Let $D$ be a probability distribution on $[0,\infty)$ with finite mean.
There is a unique $\tau_D\ge1$ satisfying
\begin{equation}
 \int_{[0,\infty)}(\tau_D-p)^+\,dD(p)=1.
 \label{eq:common-threshold-equation}
\end{equation}
Choose $g_*$ to select every positive $p<\tau_D$, reject every
$p>\tau_D$, and select any fraction of the mass at $p=\tau_D$.  Write
$a_*=a_D(g_*)$ and $w_*=w_D(g_*)$.  Then
\begin{equation}
 \frac{a_*}{w_*}=\max_g\frac{a_D(g)}{w_D(g)}=\frac1{\tau_D},
 \qquad
 \tau_D a_D(g)\le w_D(g)\quad\text{for every }g.
 \label{eq:common-density-dual}
\end{equation}
\end{lemma}

\begin{proof}
The function
$h(t)=\int(t-p)^+\,dD(p)$ is continuous and nondecreasing, satisfies
$h(0)=0$, and tends to infinity because
$h(t)\ge t-\int p\,dD(p)$.  Once positive it is strictly increasing, so
$h(t)=1$ has a unique solution.  The inequality $h(t)\le t$ gives
$\tau_D\ge1$.

For an arbitrary selector $g$,
\[
 \begin{aligned}
  \tau_D a_D(g)-\int_{(0,\infty)}p g(p)\,dD(p)
   &=\tau_D D(\{0\})
     +\int_{(0,\infty)}(\tau_D-p)g(p)\,dD(p)\\
   &\le\int_{[0,\infty)}(\tau_D-p)^+\,dD(p)=1.
 \end{aligned}
\]
This is the second inequality in \cref{eq:common-density-dual}.  The
threshold selector $g_*$ attains equality because it keeps exactly the
positive terms; selecting mass at $p=\tau_D$ changes both sides equally.
Hence $w_*=\tau_Da_*$ and $g_*$ has maximum density $1/\tau_D$.
\end{proof}

For the rest of this section, return to the finite distribution in
\cref{eq:fluid-finite-distribution}.
The positive mass not completed in the maximum-density testing block forms
the finite residual measure
\begin{equation}
 d\nu(p)=(1-g_*(p))\one_{p>0}\,dD(p),
 \qquad \nu(\{0\})=0.
 \label{eq:common-module-data}
\end{equation}
It has total mass $1-a_*$ and support in $[\tau_D,\infty)$.  Thus every
deferred type has at least as much work per completion as the testing block.

\subsection{Completion curves of finite block collections}
\label{sec:fluid-block-order}

Once testing activity has been compressed into one block, only an ordering
problem remains.  The next elementary lemma is the divisible version of the
shortest-processing-time rule.

Represent a finite block of completion mass $c_k$ and work $v_k$ per
completion by $c_k\delta_{v_k}$.  For a collection
$\mathcal M=\sum_kc_k\delta_{v_k}$ of total mass one, define
\begin{equation}
 K_{\mathcal M}(x)\defeq
 \max\left\{\sum_kz_k:
       0\le z_k\le c_k,\ \sum_kv_kz_k\le x\right\}.
 \label{eq:divisible-spt-envelope}
\end{equation}
This finite fractional-knapsack problem asks for the largest mass that the
blocks can complete using work $x$.  Its value as a function of $x$ is their
completion curve.  We also write
\begin{equation}
 \SPT(\xi)\defeq
 \frac12\iint\min\{v,w\}\,d\xi(v)d\xi(w).
 \label{eq:divisible-spt-cost}
\end{equation}

\begin{lemma}[Shortest-first order for divisible blocks]
\label{lem:divisible-spt-area}
The curve $K_{\mathcal M}$ is realized by processing the blocks in
nondecreasing order of $v_k$, and
\begin{equation}
 \Area(K_{\mathcal M})=\SPT(\mathcal M).
 \label{eq:divisible-spt-area}
\end{equation}
\end{lemma}

\begin{proof}
Consider two consecutive blocks with masses $c,c'$ and work per completion
$v\le v'$.  Putting the $v$-block first changes the area, relative to the
opposite order, by $cc'(v-v')\le0$.  Repeated neighboring exchanges sort all
blocks by nondecreasing work per completion.

In a block of mass $c_k$ and work per completion $v_k$, unfinished mass
decreases linearly.  Once the blocks are sorted, its area contribution is
\[
 v_kc_k\left(1-\sum_{h<k}c_h-\frac{c_k}{2}\right).
\]
Summing these terms and collecting each unordered pair gives
\cref{eq:divisible-spt-area}.
\end{proof}

\subsection{The optimal schedule at a fixed tested fraction}
\label{sec:fluid-fixed-q}

We can now state the common optimality theorem.  The proof compares complete
curves, not merely their final areas: after every possible amount of work,
one explicit schedule has completed at least as much mass as any competitor.

First describe the accounting state of an arbitrary schedule after it has
used work $x$.  Let $t$ be its tested mass, $b$ its mass completed in the
other mode, and
\[
 c=(c_1,\ldots,c_m)
\]
the vector of tested positive mass already processed in each class.  The
postponed vector is not missing from this state: its $i$th coordinate is
$d_it-c_i$.  An algorithm may retain a much richer internal history, but its
aggregate state necessarily satisfies
\[
 0\le t\le q,\qquad 0\le b\le1-q,\qquad 0\le c_i\le d_it.
\]
It has completed mass $d_0t+b+\sum_i c_i$ and used work
$t+vb+\sum_i p_ic_i$.  Therefore define the finite-dimensional envelope
\begin{equation}
 \mathcal C^v_{D,q}(x)\defeq
 \max\left\{
 d_0t+b+\sum_{i=1}^m c_i:
 \begin{array}{l}
 0\le t\le q,\quad 0\le b\le1-q,\\
 0\le c_i\le d_it\quad(1\le i\le m),\\
 t+vb+\sum_{i=1}^m p_ic_i\le x
 \end{array}
 \right\}.
 \label{eq:common-fluid-envelope}
\end{equation}

The candidate schedule is completely explicit.

\begin{algorithm}[H]
\caption{Shortest-first fluid schedule
$\textsc{FluidSPT}(D,q,v)$}
\label{alg:fluid-spt}
\begin{algorithmic}
\State Compute the maximum-density selector $g_*$, its completed mass $a_*$,
and its work per completion $\tau_D$.
\State Create a testing block of completion mass $qa_*$ and work per
completion $\tau_D$; running it tests mass $q$ and uses selector $g_*$.
\State For every $p_i>0$, create a deferred block of mass
$qd_i(1-g_*(p_i))$ and work per completion $p_i$.
\State Create an alternative-completion block of mass $1-q$ and work per
completion $v$.
\State Order all nonempty blocks by nondecreasing work per completion,
placing the testing block before a deferred block in case of a tie, and
execute them in that order.
\end{algorithmic}
\end{algorithm}

Equivalently, the collection of blocks created by the algorithm is
\begin{equation}
 \mathcal M_{D,q,v}\defeq
 qa_*\delta_{\tau_D}+q\nu+(1-q)\delta_v.
 \label{eq:common-envelope-items}
\end{equation}

\begin{theorem}[Pointwise optimality of the finite fluid schedule]
\label{lem:test-module-envelope}
For every complete fluid schedule $S$ that tests mass $q$,
\begin{equation}
 C_S(x)\le \mathcal C^v_{D,q}(x)
 =K_{\mathcal M_{D,q,v}}(x)
 \qquad\text{for every }x\ge0.
 \label{eq:fluid-pointwise-optimality}
\end{equation}
The schedule $\textnormal{\textsc{FluidSPT}}(D,q,v)$ attains equality for
every $x$.  Consequently it minimizes fluid total completion time and
\begin{equation}
 \Area(\mathcal C^v_{D,q})=\SPT(\mathcal M_{D,q,v}).
 \label{eq:common-envelope-area}
\end{equation}
\end{theorem}

The proof has three short parts.  First, the envelope contains the aggregate
state of every schedule.  Second, the density inequality compresses every
feasible envelope state into the blocks of
\cref{eq:common-envelope-items}.  Finally, the support property of $\nu$
makes the shortest-first block order operationally feasible.

\begin{claim}[Every schedule obeys the envelope]
\label{clm:fluid-envelope-upper}
For every complete fluid schedule $S$ and every $x\ge0$,
$C_S(x)\le\mathcal C^v_{D,q}(x)$.
\end{claim}

\begin{proof}
At work $x$, let $t,b,c_1,\ldots,c_m$ be the aggregate state described
above.  Every tested mass has type distribution $D$, so at most $d_it$ mass
of positive type $i$ has been revealed and hence $c_i\le d_it$.  The test,
alternative, and positive-processing work is respectively
$t$, $vb$, and $\sum_i p_ic_i$.  Thus the aggregate state is feasible in
\cref{eq:common-fluid-envelope}, whose objective is exactly its completed
mass.
\end{proof}

\begin{claim}[Testing activity compresses into the canonical blocks]
\label{clm:fluid-compression}
For every $x\ge0$,
$\mathcal C^v_{D,q}(x)\le K_{\mathcal M_{D,q,v}}(x)$.
\end{claim}

\begin{proof}
Fix a feasible point $(t,b,c_1,\ldots,c_m)$ in
\cref{eq:common-fluid-envelope}.  For each positive type, put
\[
 s_i=\min\{c_i,d_itg_*(p_i)\},
 \qquad r_i=c_i-s_i.
\]
The tests, tested zeros, and masses $s_i$ complete
\[
 y_{\rm test}=d_0t+\sum_i s_i
\]
using work
\[
 x_{\rm test}=t+\sum_i p_is_i.
\]
If $t>0$, apply \cref{eq:common-density-dual} to the selector
$s_i/(d_it)$, taking it to be zero when $d_i=0$.  This gives
\[
 \tau_Dy_{\rm test}\le x_{\rm test},
 \qquad y_{\rm test}\le a_*t\le a_*q.
\]
The same inequalities are trivial for $t=0$.  Hence the completed mass
$y_{\rm test}$ fits inside the canonical testing block and requires no more
work there.

Moreover,
\[
 0\le r_i\le d_it(1-g_*(p_i))
        \le qd_i(1-g_*(p_i)),
\]
so the remaining processed mass fits inside the deferred block of type $i$.
The mass $b$ fits inside the alternative block.  We have represented the
same total completed mass by items from $\mathcal M_{D,q,v}$ using no more
than the original work $x$.  The definition of
$K_{\mathcal M_{D,q,v}}$ proves the claim.
\end{proof}

\begin{claim}[The shortest-first curve is feasible]
\label{clm:fluid-realization}
The schedule $\textnormal{\textsc{FluidSPT}}(D,q,v)$ is a complete fluid
schedule and has completion curve $K_{\mathcal M_{D,q,v}}$.
\end{claim}

\begin{proof}
The alternative block can be executed wherever its value $v$ places it in
the shortest-first order.  Running any fraction of the testing block means
testing the corresponding fraction of its mass and using $g_*$ on the
revealed types.  By \cref{eq:common-module-data}, every nonempty deferred
block has processing time at least $\tau_D$.  The prescribed tie rule
therefore places the testing block before all deferred work, so each deferred
queue has been exposed before the schedule uses it.  All block capacities
are respected, and their total completion mass is
\[
 qa_*+q(1-a_*)+(1-q)=1.
\]
Thus the schedule is complete.  By
\cref{lem:divisible-spt-area}, its shortest-first execution realizes
$K_{\mathcal M_{D,q,v}}$.
\end{proof}

\begin{proof}[Proof of \cref{lem:test-module-envelope}]
The first inequality in \cref{eq:fluid-pointwise-optimality} is
\cref{clm:fluid-envelope-upper}.  Claims
\ref{clm:fluid-compression} and~\ref{clm:fluid-realization} give the equality
and show that one schedule realizes it simultaneously for every work budget.
Integrating the pointwise inequality and applying
\cref{lem:divisible-spt-area} proves
\cref{eq:common-envelope-area} and optimality.
\end{proof}

\subsection{From the fluid theorem to finite inputs}
\label{sec:fluid-specializations}

The theorem above deliberately fixes $D$, $q$, and $v$.  We next prove one
transfer theorem showing that its value is also the correct leading-order
cost for bounded shuffled inputs.  The remaining subsections then specialize
that result and perform the model-specific algebra.

In obligatory testing every job must be tested, so $q=1$ and the alternative
block is absent.  In revealing optimization with raw time $u$, set $v=u$
and minimize the area in \cref{eq:common-envelope-area} over
$q\in[0,1]$.  In blind execution, predictably chosen untouched jobs retain
distribution $D$, so blind execution of unit mass uses mean work
\[
 \mu_D=\int p\,dD(p);
\]
there we set $v=\mu_D$ and again minimize over $q$.  The resulting
model-specific values are denoted by $\PhiOT(D)$, $\PhiRO(D)$, and
$\PhiBE(D)$ in their respective sections.

For an actual vector $p$, its empirical distribution $D[p]$ is already
finitely supported.  Gridding merely reduces the number of coordinates so
that one concentration event controls them simultaneously.  The next three
subsections show that every finite competitor induces a fluid state obeying
\cref{eq:common-fluid-envelope}, up to $o(1)$ normalized error, and that a
sampled finite algorithm implements $\textnormal{\textsc{FluidSPT}}$ with
the same accuracy.  Thus the probabilistic transfer is a theorem following
the exact fluid optimization, not an assumption used inside it.
